\chapter{The covariance tensor and the new channel geometry}

\section{Two different spaces must not be confused}

The covariance-tensor approach compares whole symmetric matrices such as \(\G_1,\ldots,\G_m\). The channel approach compares invariant subspaces inside ordinary trait space after constructing \(\Sop_i\). These are related but distinct geometries:

\begin{itemize}[leftmargin=1.6em]
\item \textbf{Matrix space:} each complete \(\G_i\) or \(\Sop_i\) is one point.
\item \textbf{Trait space:} each eigenvector of a matrix is a direction; each leading eigenspace is a channel.
\end{itemize}

\begin{cautionbox}
An eigentensor is an axis of variation among matrices. A channel eigenvector is an axis of response among traits. They live in different spaces and answer different questions.
\end{cautionbox}

\section{The space of symmetric matrices}

The real symmetric \(n\times n\) matrices form a vector space \(\operatorname{Sym}(n)\) of dimension
\begin{equation}
q=\frac{n(n+1)}2.
\end{equation}
The Frobenius inner product is
\begin{equation}
\innerF{\mathbf A}{\mathbf B}
=\tr(\mathbf A^\top\mathbf B)
=\tr(\mathbf A\mathbf B)
\end{equation}
for symmetric matrices.

Define \(\svec(\mathbf A)\in\R^q\) by listing diagonal entries without scaling and one copy of every off-diagonal entry multiplied by \(\sqrt2\). For example, for a two-by-two symmetric matrix,
\begin{equation}
\mathbf A=
\begin{bmatrix}
a_{11}&a_{12}\\a_{12}&a_{22}
\end{bmatrix},
\qquad
\svec(\mathbf A)=
\begin{bmatrix}
a_{11}\\a_{22}\\\sqrt2a_{12}
\end{bmatrix}.
\end{equation}

For symmetric \(\mathbf A\) and \(\mathbf B\),
\begin{align}
\svec(\mathbf A)^\top\svec(\mathbf B)
&=\sum_iA_{ii}B_{ii}
+2\sum_{i>j}A_{ij}B_{ij}\\
&=\sum_{i,j}A_{ij}B_{ij}\\
&=\tr(\mathbf A\mathbf B)\\
&=\innerF{\mathbf A}{\mathbf B}.
\label{eq:svecIsometry}
\end{align}

\inwords The \(\sqrt2\) factor makes vectorization preserve matrix lengths and angles. Without it, off-diagonal information would be underweighted because a symmetric off-diagonal entry occurs twice in the full matrix.

\begin{connectionbox}
This is exactly why the existing covariance-tensor code multiplies covariance--covariance blocks by \(2\), cross blocks by \(\sqrt2\), and maps eigentensor off-diagonal loadings back with \(1/\sqrt2\). The code is implementing an isometry between symmetric-matrix space and ordinary Euclidean vector space.
\end{connectionbox}

\section{Construct the covariance tensor as an ordinary covariance matrix}

Let
\begin{equation}
\mathbf g_r=\svec(\G_r),
\qquad r=1,\ldots,m,
\end{equation}
and
\begin{equation}
\bar{\mathbf g}=\frac1m\sum_{r=1}^m\mathbf g_r.
\end{equation}
The covariance of matrix elements is
\begin{equation}
\mathbf C_G
=\frac1{m-1}\sum_{r=1}^m
(\mathbf g_r-\bar{\mathbf g})
(\mathbf g_r-\bar{\mathbf g})^\top.
\label{eq:matrixSpaceCovariance}
\end{equation}
This \(q\times q\) symmetric matrix is the coordinate representation of the fourth-order covariance tensor.

Diagonalize it:
\begin{equation}
\mathbf C_G\mathbf e_s=\alpha_s\mathbf e_s.
\end{equation}
Map each eigenvector back to a symmetric matrix:
\begin{equation}
\mathbf E_s=\svec^{-1}(\mathbf e_s).
\end{equation}
The matrices \(\mathbf E_s\) are eigentensors.

\section{Coordinates of each original matrix}

The coordinate of \(\G_r\) on eigentensor \(\mathbf E_s\) is
\begin{align}
c_{rs}
&=\mathbf e_s^\top\mathbf g_r\\
&=\svec(\mathbf E_s)^\top\svec(\G_r)\\
&=\innerF{\mathbf E_s}{\G_r}\\
&=\tr(\mathbf E_s\G_r).
\label{eq:eigentensorCoordinate}
\end{align}

\inwords A matrix coordinate is a dot product. Large positive and negative coordinates identify matrices that differ strongly and oppositely along an eigentensor axis.

An eigentensor can itself be diagonalized:
\begin{equation}
\mathbf E_s=\sum_{j=1}^n\eta_{sj}\mathbf e_{sj}\mathbf e_{sj}^\top.
\end{equation}
The vectors \(\mathbf e_{sj}\) identify trait combinations that contribute to the matrix-space difference captured by \(\mathbf E_s\).

\begin{geombox}
The first eigenanalysis finds directions through the cloud of whole matrices. The second eigenanalysis opens one such direction back up and asks which trait combinations create it.
\end{geombox}

\section{Three uses in the geometric snowball programme}

\subsection{Tensor on \(\G_i\)}
This reproduces the earlier question: how has the distribution of additive variance changed among populations or ecotypes? It is directly continuous with the established \textit{Senecio} analyses \cite{walter2018,aguirre2014,hine2009}.

\subsection{Tensor on \(\Hh_i\)}
This asks how local fitness curvature differs among habitats. It can identify recurrent combinations of traits in which selection surfaces diverge.

\subsection{Tensor on \(\Sop_i\)}
This asks how the complete joint response operators differ. It is useful for describing operator-level divergence across many populations. Principal-angle analysis is then applied to the leading eigenspaces of each \(\Sop_i\) to quantify channel mismatch directly.

\begin{figure}[H]
\centering
\includegraphics[width=\textwidth]{figures/pedagogy_tensor_bridge.pdf}
\caption{Bridge from the established covariance-tensor framework to response-operator and channel geometry. Eigentensors describe axes of whole-matrix divergence; principal angles describe divergence among response subspaces.}
\label{fig:tensorBridge}
\end{figure}

\section{A new synthesis statistic}

Across multiple populations, one can separate three levels:
\begin{enumerate}[leftmargin=1.6em]
\item \(\alpha_s\): how much matrices differ along eigentensor \(s\);
\item \(c_{rs}\): which populations occupy opposite ends of that matrix-space axis;
\item \(\psi_{ij,1},\ldots,\psi_{ij,k}\): the principal angles measuring how much the biologically chosen response subspaces of populations \(i\) and \(j\) differ.
\end{enumerate}

This makes it possible to ask whether the major axes of operator divergence identified by a tensor are the same changes that generate large channel angles and hybrid load.
